\chapter{Fundamentação Teórica}
\label{chap:fundamentacao}

\begin{resumocapitulo}
Este capítulo constrói, do concreto ao abstrato, todo o vocabulário técnico
de que a dissertação precisa: o que são dados tabulares e por que eles
resistem às arquiteturas que revolucionaram visão e linguagem
(\S\ref{sec:fund-tabular}); como uma única árvore de decisão aprende e por
que \emph{ensembles} de árvores com \emph{boosting} se tornaram o padrão da
área (\S\ref{sec:fund-gbdt}); as três linhagens da resposta do aprendizado
profundo, do MLP aos Transformers e às redes de Kolmogorov--Arnold
(\S\ref{sec:fund-dl}); os \emph{foundation models} tabulares e o mecanismo
de aprendizado em contexto que dispensa treinamento na tarefa
(\S\ref{sec:fund-fm}); a destilação de conhecimento
(\S\ref{sec:fund-destilacao}); o ferramental da regressão distribucional
--- perda \emph{pinball}, CRPS, calibração --- que sustenta a contribuição
do Capítulo~7 (\S\ref{sec:fund-distribucional}); e os testes estatísticos
com que o Capítulo~5 compara os modelos (\S\ref{sec:fund-estatistica}). A
profundidade matemática cresce com a proximidade das contribuições: o
\emph{boosting} é apresentado verbalmente, a atenção com suas equações
completas, e a regressão distribucional e a inferência estatística com o
detalhe necessário para ler os resultados sem consultar as fontes.
\end{resumocapitulo}

\publicacaoassociada{As formulações de CRPS, PICP, perda \emph{pinball} e
aprendizado em contexto das Seções~\ref{sec:fund-fm}
e~\ref{sec:fund-distribucional} seguem as glosas auditadas dos dois artigos
que acompanham esta dissertação (\texttt{latex/ai1} e \texttt{latex/ai2} no
repositório), aqui traduzidas e expandidas didaticamente.}

\paragraph{Roteiro.} O leitor já familiarizado com \emph{gradient boosting}
e redes neurais pode saltar diretamente para a
Seção~\ref{sec:fund-fm}, onde começa o material menos difundido
(aprendizado em contexto e cabeças distribucionais); quem conhece os
\emph{foundation models} tabulares, para as
Seções~\ref{sec:fund-distribucional} e~\ref{sec:fund-estatistica}, que são
pré-requisito direto dos Capítulos~5 e~7. As demais seções foram escritas
para que um leitor com formação geral em computação acompanhe toda a
dissertação sem fontes externas.

\section{Dados tabulares e aprendizagem supervisionada}
\label{sec:fund-tabular}

Comecemos pelo objeto. Um \emph{conjunto de dados tabular} é uma tabela em
que cada linha é uma observação e cada coluna um atributo --- o formato de
um extrato bancário, de um prontuário ou de um cadastro de clientes. A
Tabela~\ref{tab:adult-exemplo} ilustra o formato com linhas no estilo do
conjunto \texttt{adult}, um clássico da área em que se deseja prever, a
partir de atributos sociodemográficos de um censo, se a renda anual de uma
pessoa excede 50 mil dólares \textbf{[REF-NOVA: Kohavi 1996, conjunto
Adult/Census Income do repositório UCI]}.

\begin{table}[htb]
\centering
\caption{Fragmento \emph{ilustrativo} de um conjunto tabular no estilo do
\texttt{adult}. As linhas são fictícias, porém plausíveis; servem apenas
para fixar o formato: atributos numéricos e categóricos convivem na mesma
tabela, há valores ausentes (``?'') e a última coluna é o alvo a prever.}
\label{tab:adult-exemplo}
\begin{tabular}{rllrl}
\toprule
idade & escolaridade & ocupação & horas/semana & renda $>$ 50K? \\
\midrule
39 & superior        & gerência        & 45 & sim \\
23 & ensino médio    & serviços        & 40 & não \\
51 & pós-graduação   & especialista    & 50 & sim \\
34 & ensino médio    & ?               & 38 & não \\
28 & superior        & técnico         & 44 & não \\
\bottomrule
\end{tabular}
\end{table}

\paragraph{A heterogeneidade é constitutiva.} A Tabela~\ref{tab:adult-exemplo}
já exibe tudo o que torna o formato difícil: \texttt{idade} e
\texttt{horas/semana} são números em escalas distintas;
\texttt{escolaridade} é uma categoria com ordem natural;
\texttt{ocupação} é uma categoria \emph{nominal}, sem ordem nem distância
entre valores, e de cardinalidade arbitrária; e o ``?'' é um valor ausente
cuja ausência pode, ela própria, carregar informação. Uma mesma tabela
mistura, portanto, variáveis contínuas, contagens, categorias e lacunas ---
e cada coluna é um eixo semântico independente: trocar duas colunas de
lugar não altera o problema, mas trocar dois pixels de uma imagem, sim.

\paragraph{Por que as revoluções de visão e linguagem não transferiram?}
Redes convolucionais exploram a estrutura espacial das imagens (pixels
vizinhos são correlacionados; padrões se repetem em qualquer posição);
Transformers de linguagem exploram a estrutura sequencial do texto. Dados
tabulares não oferecem nenhuma invariância universal análoga que uma
arquitetura possa explorar \emph{por construção} --- não há noção de
vizinhança entre \texttt{idade} e \texttt{ocupação}. Essa ausência de
estrutura compartilhada é a explicação usual para um fato empírico
persistente: por mais de uma década, o estado da arte em tabelas pertenceu
a outra família de modelos, os \emph{ensembles} de árvores de decisão
\citep{grinsztajn2022,mcelfresh2023}. As Seções~\ref{sec:fund-gbdt}
e~\ref{sec:fund-dl} desenvolvem os dois lados dessa disputa.

\paragraph{O problema da aprendizagem supervisionada.} Formalizemos o
mínimo necessário. Dado um conjunto de treino
$D = \{(\mathbf{x}_i, y_i)\}_{i=1}^{n}$ de pares observação--alvo,
assumidos amostras independentes de uma distribuição desconhecida, o
objetivo é aprender uma função $f$ que prediga bem o alvo de observações
\emph{novas} --- isto é, que minimize o erro esperado sob a distribuição
geradora, não o erro no próprio conjunto de treino. A distinção é o coração
do problema: um modelo suficientemente flexível pode zerar o erro de treino
simplesmente memorizando $D$ (\emph{sobreajuste}), sem generalizar. Por
isso toda avaliação honesta mede o desempenho em dados que o modelo --- e,
igualmente importante, o processo de \emph{escolha} do modelo --- jamais
viu; o Capítulo~4 traduz esse princípio em protocolo (conjunto de teste
separado antes de tudo, validação cruzada aninhada para a seleção de
hiperparâmetros).

\paragraph{As três tarefas.} Este trabalho considera os três tipos
canônicos de tarefa supervisionada sobre tabelas: \emph{classificação
binária} (prever uma de duas classes, como na
Tabela~\ref{tab:adult-exemplo}), \emph{classificação multiclasse} (prever
uma entre $K$ classes) e \emph{regressão} (prever um valor contínuo). A
regressão admite duas leituras que esta dissertação distingue com cuidado:
a leitura \emph{pontual}, em que o modelo devolve um único número (uma
estimativa da média condicional), e a leitura \emph{distribucional}, em que
o modelo devolve uma distribuição preditiva completa --- quantis,
intervalos, densidades --- da qual se exige calibração. A distinção parece
sutil agora, mas é o eixo da contribuição do Capítulo~7, e a
Seção~\ref{sec:fund-distribucional} lhe dedica o ferramental próprio.

\section{Das árvores de decisão ao \emph{gradient boosting}}
\label{sec:fund-gbdt}

\subsection{Uma árvore, um exemplo guiado}

Uma \emph{árvore de decisão} prediz por meio de uma sequência de perguntas
sobre os atributos, cada uma com resposta binária, até chegar a uma folha
que contém a predição \textbf{[REF-NOVA: Breiman, Friedman, Olshen \& Stone
1984, CART]}. A Figura~\ref{fig:arvore-exemplo} mostra uma árvore pequena,
de profundidade dois, treinada (ilustrativamente) sobre dados como os da
Tabela~\ref{tab:adult-exemplo}.

\begin{figure}[htb]
\centering
\begin{tikzpicture}[
  level distance=17mm,
  level 1/.style={sibling distance=62mm},
  level 2/.style={sibling distance=31mm},
  interno/.style={draw, rounded corners=2pt, align=center, font=\small,
                  inner sep=4pt},
  folha/.style={draw, rounded corners=2pt, align=center, font=\small,
                fill=black!6, inner sep=4pt},
  rotulo/.style={font=\footnotesize\itshape, midway},
  edge from parent/.style={draw, -{Stealth[length=2mm]}}]
\node[interno] {escolaridade $\geq$ superior?}
  child { node[interno] {horas/semana $>$ 45?}
    child { node[folha] {$\hat{p} = 0{,}06$}
      edge from parent node[rotulo, left=1pt] {não} }
    child { node[folha] {$\hat{p} = 0{,}22$}
      edge from parent node[rotulo, right=1pt] {sim} }
    edge from parent node[rotulo, left=3pt] {não} }
  child { node[interno] {idade $>$ 30?}
    child { node[folha] {$\hat{p} = 0{,}29$}
      edge from parent node[rotulo, left=1pt] {não} }
    child { node[folha] {$\hat{p} = 0{,}61$}
      edge from parent node[rotulo, right=1pt] {sim} }
    edge from parent node[rotulo, right=3pt] {sim} };
\end{tikzpicture}
\caption{\textbf{Uma árvore de decisão responde por partições sucessivas do
espaço de atributos.} Árvore ilustrativa de profundidade 2 sobre o exemplo
da Tabela~\ref{tab:adult-exemplo}; em cada folha, $\hat{p}$ é a fração de
exemplos de treino daquela região com renda $>$ 50K (valores fictícios,
para ilustração). Para classificar uma linha nova, percorre-se o caminho
ditado pelas respostas e lê-se a predição na folha alcançada.}
\label{fig:arvore-exemplo}
\end{figure}

Siga a terceira linha da Tabela~\ref{tab:adult-exemplo} pela árvore: a
pessoa tem pós-graduação (escolaridade $\geq$ superior: \emph{sim}, ramo
direito) e 51 anos (idade $>$ 30: \emph{sim}), caindo na folha com
$\hat{p} = 0{,}61$ --- a árvore estima 61\% de probabilidade de renda alta,
e, com limiar de 50\%, prediz ``sim''. O treinamento constrói essas
perguntas de cima para baixo, escolhendo em cada nó, por busca gulosa, o
atributo e o ponto de corte que melhor separam as classes (segundo um
critério de impureza como entropia ou índice de Gini); em regressão, as
folhas guardam médias e o critério é a redução de variância.

Duas propriedades dessa construção merecem registro desde já, porque
reaparecerão como explicação de resultados empíricos. Primeiro, as
partições são \emph{alinhadas aos eixos}: cada pergunta envolve um único
atributo, o que casa com a natureza tabular em que cada coluna tem
semântica própria. Segundo, a árvore é \emph{invariante a transformações
monótonas} de cada atributo --- perguntar ``idade $>$ 30'' é equivalente a
perguntar ``$\log(\text{idade}) > \log 30$'' --- de modo que normalizar
colunas, etapa obrigatória para redes neurais, é irrelevante para árvores.

\subsection{Por que uma árvore só não basta?}

Uma árvore única enfrenta um dilema. Se for rasa, é estável porém grosseira:
o exemplo da Figura~\ref{fig:arvore-exemplo} reduz toda a população a
quatro grupos. Se for profunda, torna-se expressiva porém instável:
memoriza particularidades da amostra de treino, e uma pequena perturbação
nos dados produz uma árvore muito diferente. Na linguagem clássica da
decomposição \emph{viés--variância}, árvores rasas sofrem de viés alto e
árvores profundas, de variância alta \textbf{[REF-NOVA: Hastie, Tibshirani
\& Friedman 2009, \emph{The Elements of Statistical Learning}]}. Os
\emph{ensembles} --- comitês de árvores --- atacam cada chifre do dilema
por um caminho distinto.

\paragraph{Caminho 1: reduzir a variância em paralelo.} O \emph{bagging} e
seu refinamento mais célebre, as florestas aleatórias, treinam muitas
árvores profundas em réplicas perturbadas dos dados (reamostragem com
reposição e sorteio de atributos por nó) e tiram a média das predições: os
erros idiossincráticos de cada árvore, aproximadamente independentes,
cancelam-se na média, e a variância do comitê despenca sem elevar o viés
\textbf{[REF-NOVA: Breiman 2001, \emph{Random Forests}]}.

\paragraph{Caminho 2: reduzir o viés em sequência.} O \emph{boosting}
inverte a lógica: parte de árvores \emph{rasas} (viés alto, variância
baixa) e as treina em sequência, cada nova árvore concentrada em corrigir
os erros que o comitê acumulado ainda comete. A predição final é a soma
ponderada de todas as correções. A cada rodada, os exemplos mal preditos
ganham protagonismo, e o comitê refina exatamente as regiões em que ainda
erra --- o viés cai passo a passo, enquanto a rasura de cada árvore e um
passo de aprendizado pequeno seguram a
variância.\footnote{Tecnicamente, o \emph{gradient boosting} formaliza essa
intuição como uma descida de gradiente no espaço de funções: a cada rodada
$m$, ajusta-se uma árvore $h_m$ ao gradiente negativo da função de perda
avaliada nas predições correntes (no caso do erro quadrático, aos resíduos
$y_i - F_{m-1}(x_i)$), e atualiza-se $F_m = F_{m-1} + \nu\, h_m$ com taxa
de aprendizado $\nu$ pequena \textbf{[REF-NOVA: Friedman 2001,
\emph{Greedy Function Approximation: A Gradient Boosting Machine}]}. O
corpo do texto permanece verbal de propósito: nenhum resultado desta
dissertação exige manipular essa formulação.}

\subsection{XGBoost, LightGBM e CatBoost: uma diferença concreta cada}

Três implementações modernas de \emph{gradient boosting} sobre árvores
(GBDTs, de \emph{gradient-boosted decision trees}) dominam a prática, e o
benchmark do Capítulo~5 avalia as três. Para os fins desta dissertação,
basta reter uma diferença concreta de cada uma. O \textbf{XGBoost}
\citep{chen2016} consolidou o paradigma com uma formulação regularizada do
objetivo --- penalidades explícitas sobre o número de folhas e a magnitude
dos pesos --- que disciplina o crescimento de cada árvore. O
\textbf{LightGBM} \citep{ke2017} muda a ordem de crescimento: em vez de
expandir a árvore nível a nível, expande sempre a \emph{folha} de maior
ganho (crescimento folha-a-folha), com histogramas de atributos que tornam
o treino substancialmente mais rápido em dados grandes. O
\textbf{CatBoost} \citep{prokhorenkova2018} distingue-se pelo tratamento
nativo de variáveis categóricas via \emph{ordered boosting}: as estatísticas
de alvo usadas para codificar categorias são computadas apenas sobre
exemplos ``anteriores'' numa permutação aleatória, evitando o vazamento de
alvo das codificações ingênuas.

\paragraph{Os botões de ajuste.} Na prática, o comportamento de um GBDT é
governado por um punhado de hiperparâmetros com papéis interpretáveis à luz
do dilema viés--variância: a \emph{profundidade} das árvores controla a
ordem das interações que cada correção pode capturar; o \emph{número de
árvores} e a \emph{taxa de aprendizado} formam um par acoplado (passos
menores exigem mais passos, mas regularizam); e as frações de
\emph{subamostragem} de linhas e colunas por árvore injetam a
descorrelação que o \emph{bagging} ensinou. Nenhuma configuração domina
todos os conjuntos de dados --- razão pela qual qualquer comparação justa
entre modelos precisa otimizar os hiperparâmetros de cada um sob o mesmo
orçamento, como especifica o Capítulo~4.

\paragraph{Por que o \emph{boosting} domina dados tabulares?} Reunindo os
fios: as partições alinhadas aos eixos capturam não-linearidades e
interações por atributo sem supor qualquer geometria entre colunas; a
invariância a transformações monótonas dispensa normalização e absorve
escalas heterogêneas; a seleção gulosa de cortes ignora com naturalidade
atributos irrelevantes; e o \emph{boosting} compõe essas virtudes num
redutor sistemático de viés. O resultado é um desempenho consistente, com
pouco ajuste, exatamente no regime de dados heterogêneos de tamanho pequeno
a médio --- diagnóstico corroborado por avaliações independentes em larga
escala \citep{grinsztajn2022,mcelfresh2023}. Por isso, nesta dissertação os
GBDTs cumprem o papel de \emph{hipótese nula}: toda proposta nova em dados
tabulares deve ser medida contra eles.

\section{Aprendizado profundo para dados tabulares}
\label{sec:fund-dl}

A resposta do aprendizado profundo ao domínio dos GBDTs organizou-se em
três linhagens: modernizar o MLP, importar a atenção dos Transformers e,
mais recentemente, repensar o próprio neurônio com as redes de
Kolmogorov--Arnold. Esta seção as percorre em ordem.

\subsection{O perceptron multicamadas}

O \emph{perceptron multicamadas} (MLP) é a rede neural mais simples:
camadas de unidades que combinam linearmente as entradas e aplicam uma
não-linearidade. Com entrada $\mathbf{h}^{(0)} = \mathbf{x}$, cada camada
$l$ computa
\begin{equation}
  \mathbf{h}^{(l)} \;=\; \phi\!\left(\mathbf{W}^{(l)} \mathbf{h}^{(l-1)}
  + \mathbf{b}^{(l)}\right),
  \label{eq:mlp-camada}
\end{equation}
em que $\mathbf{W}^{(l)}$ e $\mathbf{b}^{(l)}$ são parâmetros aprendidos e
$\phi$ é uma função de ativação fixa (tipicamente a ReLU,
$\phi(z) = \max(0, z)$), e a saída da última camada é lida como predição,
\begin{equation}
  f(\mathbf{x}) \;=\; \mathbf{W}^{(L)} \mathbf{h}^{(L-1)} + \mathbf{b}^{(L)},
  \label{eq:mlp-saida}
\end{equation}
seguida de uma \emph{softmax} em classificação. Os parâmetros são ajustados
por descida de gradiente sobre uma função de perda, com gradientes
computados por retropropagação \textbf{[REF-NOVA: Rumelhart, Hinton \&
Williams 1986, retropropagação]}. Em uma frase: com largura suficiente, um
MLP pode aproximar qualquer função contínua \textbf{[REF-NOVA: Hornik,
Stinchcombe \& White 1989, aproximação universal]} --- a questão em dados
tabulares nunca foi expressividade, e sim se o treinamento encontra, com
dados finitos e heterogêneos, soluções tão boas quanto as dos GBDTs.

\paragraph{O gargalo do encoding categórico.} Antes de qualquer camada, há
uma decisão incômoda: como entregar \texttt{ocupação} a uma rede que só
consome números? A codificação \emph{one-hot} --- uma coluna binária por
categoria --- explode a dimensionalidade quando a cardinalidade é alta: no
benchmark desta dissertação, o conjunto \texttt{amazon\_employee} atingiria
cerca de 6{,}9 mil colunas sem um teto de cardinalidade, penalizando
especialmente os modelos de atenção, cujo custo cresce com o número de
colunas. A alternativa padrão são \emph{embeddings} aprendidos: cada
categoria é mapeada a um vetor denso de baixa dimensão, treinado junto com
a rede --- solução que todas as arquiteturas a seguir adotam de alguma
forma.

\paragraph{A receita importa: RealMLP e TabM.} Duas propostas recentes
mostram que boa parte da desvantagem histórica do MLP era de \emph{receita},
não de arquitetura. O \textbf{RealMLP} \citep{realmlp2024} demonstra que um
MLP com um conjunto cuidadoso de melhorias --- \emph{embeddings} numéricos,
agendamento de taxa de aprendizado, regularização e pré-processamento
afinados --- compete com GBDTs ajustados. O \textbf{TabM} \citep{tabm2025}
ataca outro flanco: obtém o efeito estabilizador de um \emph{ensemble} de
MLPs com custo próximo ao de um modelo único, parametrizando eficientemente
múltiplas ``sub-redes'' que compartilham a maior parte dos pesos.

\paragraph{Por que insistir em redes para tabelas?} Se os GBDTs dominam,
por que a área continua investindo em redes? Por três promessas que
árvores não fazem. Primeiro, redes aprendem \emph{representações}:
\emph{embeddings} de categorias e camadas intermediárias reutilizáveis,
porta de entrada para transferência entre tarefas e para dados mistos
(tabela junto a texto ou imagem). Segundo, redes são treináveis por
gradiente de \emph{qualquer} perda diferenciável, o que permite cabeças de
saída flexíveis --- multi-alvo, multi-quantil, distribucionais --- com
mudança mínima de arquitetura; essa plasticidade é explorada diretamente
pelos alunos de destilação do Capítulo~7. Terceiro, e decisivo para esta
dissertação: toda a família de \emph{foundation models} da
Seção~\ref{sec:fund-fm}, que redefine o custo de entrada em novas tarefas,
só é possível sobre redes.

\subsection{Atenção e Transformers}

A segunda linhagem importa o mecanismo que revolucionou o processamento de
linguagem: a \emph{atenção} \textbf{[REF-NOVA: Vaswani et al. 2017,
\emph{Attention Is All You Need}]}. Aqui a profundidade matemática se
justifica: metade dos modelos do benchmark e todos os \emph{foundation
models} da Seção~\ref{sec:fund-fm} são construídos sobre essas equações.

Dada uma sequência de $n$ vetores de entrada empilhados em
$\mathbf{X} \in \mathbb{R}^{n \times d}$, a atenção projeta cada vetor em
três papéis --- \emph{consulta} (\emph{query}), \emph{chave} (\emph{key}) e
\emph{valor} (\emph{value}):
\begin{equation}
  \mathbf{Q} = \mathbf{X}\mathbf{W}_Q, \qquad
  \mathbf{K} = \mathbf{X}\mathbf{W}_K, \qquad
  \mathbf{V} = \mathbf{X}\mathbf{W}_V,
  \label{eq:qkv}
\end{equation}
com matrizes de projeção aprendidas, e produz, para cada posição, uma média
dos valores ponderada pela afinidade entre sua consulta e todas as chaves:
\begin{equation}
  \operatorname{Atenção}(\mathbf{Q}, \mathbf{K}, \mathbf{V})
  \;=\; \operatorname{softmax}\!\left(
    \frac{\mathbf{Q}\mathbf{K}^{\top}}{\sqrt{d_k}}
  \right)\mathbf{V},
  \label{eq:atencao}
\end{equation}
em que $d_k$ é a dimensão das chaves e a divisão por $\sqrt{d_k}$ estabiliza
as magnitudes antes da \emph{softmax}. A intuição: cada elemento da
sequência ``pergunta'' (via $\mathbf{Q}$) quais outros elementos lhe são
relevantes (via $\mathbf{K}$) e agrega a informação deles (via $\mathbf{V}$)
na proporção dessa relevância --- pesos que são \emph{computados dos
dados}, não fixados pela arquitetura. A atenção \emph{multi-cabeça} executa
$h$ atenções em paralelo com projeções independentes e concatena os
resultados,
\begin{equation}
  \operatorname{MultiHead}(\mathbf{X}) \;=\;
  \operatorname{concat}(\text{cabeça}_1, \ldots, \text{cabeça}_h)\,
  \mathbf{W}_O,
  \label{eq:multihead}
\end{equation}
permitindo que cabeças distintas capturem padrões de interação distintos. Um
\emph{bloco Transformer} empilha atenção multi-cabeça e um MLP
posição-a-posição, com conexões residuais e normalização de camada; a
arquitetura completa original empilha esses blocos em um \emph{encoder}
(que processa a entrada inteira de uma vez, com atenção irrestrita) e um
\emph{decoder} (que gera texto token a token, com atenção causal).

\paragraph{Por que dados tabulares usam só o \emph{encoder}?} Porque não há
nada a \emph{gerar} sequencialmente: a predição sobre uma linha é um único
rótulo ou valor, não uma sequência. Os Transformers tabulares tratam as
\emph{colunas} de uma linha como a ``sequência'' de entrada --- cada
atributo vira um token --- e usam apenas a pilha de \emph{encoder}: a
atenção entre tokens-atributo aprende interações entre colunas, e a
predição é lida de um token especial agregador. O \emph{decoder} e sua
máscara causal, feitos para a ordem temporal do texto, não têm contraparte
numa linha de tabela.

\paragraph{As arquiteturas de atenção do benchmark.} O \textbf{TabNet}
\citep{arik2021} foi pioneiro em levar atenção a tabelas, combinando-a com
seleção \emph{esparsa} de atributos: a cada passo de processamento, uma
máscara aprendida elege poucos atributos relevantes, o que confere ao
modelo interpretabilidade nativa. O \textbf{FT-Transformer}
\citep{gorishniy2021} é a destilação mais limpa da ideia: um
\emph{Feature Tokenizer} converte cada atributo (numérico ou categórico) em
um token-\emph{embedding} e um Transformer padrão processa a linha; é a
referência arquitetural da linhagem. O \textbf{SAINT} \citep{somepalli2021}
acrescenta à atenção entre colunas uma atenção \emph{entre linhas}
(\emph{intersample}): ao classificar uma linha, o modelo pode consultar
outras linhas do lote --- um antecipo, em miniatura, da ideia de usar
exemplos como contexto que a Seção~\ref{sec:fund-fm} leva ao extremo. O
\textbf{STab} \citep{stab-ref} completa a linhagem no benchmark
[VERIFICAR: descrição do STab --- o rascunho-fonte (\texttt{dissertation/}\allowbreak\texttt{cap2-fundamentacao.md}) o descreve como competição estocástica local
(LWTA) de Voskou et al., mas a entrada \texttt{stab-ref} do
\texttt{references.bib} aponta para Hajiramezanali et al. 2022,
aprendizado autossupervisionado para tabelas; conferir qual STab foi
implementado e alinhar descrição e referência].

\subsection{Redes de Kolmogorov--Arnold}

A terceira linhagem, mais recente, muda o lugar da não-linearidade. O
teorema da representação de Kolmogorov--Arnold garante que toda função
contínua de várias variáveis pode ser escrita como composição de somas de
funções contínuas de \emph{uma} variável \textbf{[REF-NOVA: Kolmogorov 1957
e Arnold, teorema da representação]}. As \emph{Kolmogorov--Arnold Networks}
(KANs) \citep{kan2024} tomam o teorema como inspiração arquitetural e
invertem o desenho do MLP: em vez de ativações fixas nos nós e pesos
escalares nas arestas (Eq.~\ref{eq:mlp-camada}), colocam \emph{funções
univariadas aprendíveis} nas arestas --- parametrizadas por B-splines ---
e meras somas nos nós. Cada conexão deixa de ser um número e passa a ser
uma pequena curva ajustável. Variantes trocam a base das funções de aresta
(polinômios de Chebyshev no ChebyKAN, funções de base radial no FastKAN), e
o \textbf{TabKAN} \citep{tabkan2025} organiza essas variantes num arcabouço
voltado a dados tabulares.

\paragraph{Uma recepção controversa.} A promessa das KANs --- mais
expressividade por parâmetro e interpretabilidade das curvas aprendidas ---
encontrou ceticismo empírico: sob comparação com paridade de parâmetros e
de FLOPs, o MLP vence as KANs na maioria das tarefas de aprendizado de
máquina \citep{kancritical2024}, e avaliações independentes em tabelas
chegaram a quadro semelhante \citep{poeta2024}. Nenhuma dessas avaliações,
porém, inclui GBDTs e \emph{foundation models} no mesmo protocolo --- a
lacuna que o Capítulo~5 preenche, adjudicando a controvérsia sob protocolo
neutro.

\section{\emph{Foundation models} tabulares e aprendizado em contexto}
\label{sec:fund-fm}

Tudo o que foi visto até aqui compartilha uma premissa tão básica que passa
despercebida: para cada novo conjunto de dados, treina-se um modelo novo.
Os \emph{foundation models} tabulares abandonam essa premissa, e é
essencial entender o mecanismo --- não apenas o efeito --- porque tanto o
custo de inferência medido no Capítulo~5 quanto a destilação do Capítulo~7
decorrem diretamente dele.

\paragraph{O mecanismo: treinar virou condicionar.} Um modelo de
\emph{aprendizado em contexto} (\emph{in-context learning}, ICL) é uma rede
pré-treinada que condiciona no conjunto de treino como se fosse um
\emph{prompt}, no momento da predição, sem nenhuma atualização de gradiente
na tarefa-alvo. Concretamente: apresenta-se ao Transformer, numa única
entrada, o conjunto de treino inteiro --- pares $(\mathbf{x}_i, y_i)$ ---
seguido da linha de consulta $\mathbf{x}_{\text{novo}}$; um único passe
direto (\emph{forward pass}) produz a predição, com os pesos da rede
congelados. ``Treinar'' na tarefa reduz-se a montar o contexto; todo o
aprendizado propriamente dito aconteceu antes, no pré-treinamento. A
Figura~\ref{fig:icl-mecanismo} esquematiza as duas fases.

\begin{figure}[htb]
\centering
\begin{tikzpicture}[
  font=\small,
  caixa/.style={draw, rounded corners=2pt, align=center, inner sep=5pt,
                minimum height=9mm},
  destaque/.style={caixa, fill=black!6},
  rotulo/.style={font=\footnotesize\itshape, midway, fill=white,
                 inner sep=1pt},
  fase/.style={font=\footnotesize\scshape},
  seta/.style={-{Stealth[length=2.2mm]}, semithick},
  node distance=7mm and 11mm]

% Fase 1: pré-treinamento
\node[caixa] (prior) {Prior sintético\\(modelos causais\\estruturais)};
\node[caixa, right=of prior] (tarefas) {Milhões de tarefas\\sintéticas\\
  $(X, y)$ amostradas};
\node[destaque, right=of tarefas] (modelo) {Transformer\\(pesos $\theta$)};
\node[fase, above=2mm of prior.north west, anchor=south west]
  {Fase 1 --- pré-treinamento (uma única vez, \emph{offline})};

% Fase 2: inferência
\node[caixa, below=14mm of prior] (contexto)
  {Contexto:\\treino $(X_{\text{tr}}, y_{\text{tr}})$};
\node[caixa, right=of contexto] (consulta)
  {Consulta:\\$x_{\text{novo}}$};
\node[destaque, right=of consulta] (forward)
  {Um passe direto\\($\theta$ congelado)};
\node[caixa, right=of forward] (pred)
  {Distribuição\\preditiva\\$p(y \mid x_{\text{novo}}, D)$};
\node[fase, above=2mm of contexto.north west, anchor=south west]
  {Fase 2 --- inferência em cada tarefa nova (sem gradiente)};

\draw[seta] (prior) -- (tarefas);
\draw[seta] (tarefas) -- node[rotulo, above=2pt] {treino} (modelo);
\draw[seta] (modelo.south) -- ++(0,-4mm) -| node[rotulo, near start, above]
  {$\theta$ fixado} (forward.north);
\draw[seta] (contexto) -- (consulta);
\draw[seta] (consulta) -- (forward);
\draw[seta] (forward) -- (pred);
\end{tikzpicture}
\caption{\textbf{No aprendizado em contexto, o conjunto de treino entra
como entrada do modelo, não como alvo de otimização.} Na fase 1, o
Transformer é pré-treinado uma única vez sobre milhões de tarefas
sintéticas amostradas de um prior. Na fase 2, diante de uma tarefa real, o
conjunto de treino e a linha de consulta são concatenados na entrada e um
único passe direto, sem qualquer atualização de gradiente, produz a
distribuição preditiva.}
\label{fig:icl-mecanismo}
\end{figure}

\paragraph{De onde vem a capacidade de ``aprender'' sem gradiente?} Do
pré-treinamento. As \emph{prior-fitted networks} (PFNs) são pré-treinadas
em milhões de pequenos conjuntos de dados \emph{sintéticos}, amostrados de
um prior sobre mecanismos geradores --- tipicamente modelos causais
estruturais (SCMs): grafos causais aleatórios cujas variáveis se relacionam
por funções e ruídos sorteados \textbf{[REF-NOVA: Müller et al. 2022,
\emph{Transformers Can Do Bayesian Inference} / PFNs]}. Em cada tarefa
sintética, a rede vê um contexto rotulado e é treinada a prever os rótulos
retidos; ao fazê-lo em escala, aprende um algoritmo genérico de inferência:
a saída do passe direto aproxima a \emph{distribuição preditiva posterior}
sob o prior sintético, dado o contexto \textbf{[REF-NOVA: Nagler 2023,
fundamentos estatísticos das PFNs]}. O \textbf{TabPFN} \citep{hollmann2025}
e suas versões v2/v2.5 \citep{priorlabs2025} estabeleceram o paradigma na
prática, com um limite nativo de contexto (50 mil linhas na v2.5) herdado
do custo quadrático da atenção.

\paragraph{Cabeça distribucional vs.\ cabeça escalar.} Um detalhe
arquitetural do TabPFN é central ao Capítulo~7: em regressão, sua camada de
saída não devolve um número, mas uma \emph{bar distribution} --- um
histograma discretizado sobre o intervalo do alvo, do qual se leem quantis
arbitrários da distribuição preditiva. O modelo entrega, nativamente,
regressão \emph{distribucional} (Seção~\ref{sec:fund-tabular}). Em
contraste, o \textbf{TabFM} \citep{tabfm2026,kong2026} --- implementação
independente do paradigma ICL, com arquitetura híbrida que combina atenção
entre colunas, compressão de cada linha em tokens e um Transformer de ICL,
e \emph{checkpoints} separados por tarefa --- decodifica em regressão um
único escalar: é um professor exclusivamente \emph{pontual}. Essa
assimetria de cabeças de saída, e não qualquer diferença de acurácia média,
é o que tornará o TabPFN o professor natural da destilação distribucional
do Capítulo~7. Registre-se ainda que ambos os modelos carregam limites
arquiteturais herdados do pré-treinamento --- o TabFM, por exemplo, admite
no máximo 10 classes em classificação, e até a escrita desta dissertação
está documentado em relatório técnico, sem artigo revisado por pares
\citep{tabfm2026,kong2026} --- limites que se tornam células ausentes do
benchmark e exigem a política estatística declarada no Capítulo~4.

\paragraph{O \emph{trade-off} constitutivo.} O mecanismo da
Figura~\ref{fig:icl-mecanismo} tem um preço estrutural: \emph{o conjunto de
treino viaja com o modelo a cada predição}. O custo de treinamento é
próximo de zero --- montar o contexto ---, mas cada passe de inferência
reprocessa o contexto inteiro, o que coloca o custo por predição ordens de
magnitude acima do de um GBDT, cuja inferência percorre algumas árvores. Em
cenários de inferência intensiva, essa conta domina qualquer vantagem de
acurácia --- o Capítulo~5 a quantifica e o Capítulo~7 propõe uma saída.

\section{Destilação de conhecimento}
\label{sec:fund-destilacao}

Se um modelo caro é bom demais para descartar e caro demais para servir, a
resposta clássica é a \emph{destilação de conhecimento}: treinar um modelo
rápido (o \emph{aluno}) para imitar as saídas de um modelo caro (o
\emph{professor}). A ideia nasce como \emph{compressão de modelos}
\citep{bucila2006} --- um \emph{ensemble} volumoso rotula uma grande massa
de dados não rotulados, e uma rede pequena é treinada nesses rótulos ---
e ganha sua formulação moderna com \citet{hinton2015}.

\paragraph{\emph{Soft targets}: por que imitar é mais fácil que aprender.}
A observação central de \citet{hinton2015} é que as \emph{probabilidades}
que o professor atribui às classes erradas carregam informação que o rótulo
verdadeiro não tem: ao classificar a quarta linha da
Tabela~\ref{tab:adult-exemplo}, um professor que responde ``não, com 85\%''
ensina mais do que o rótulo seco ``não'' --- codifica o quão ambíguo é o
caso. Para expor essa estrutura, amplifica-se a entropia da saída com uma
\emph{temperatura} $T$ na \emph{softmax},
\begin{equation}
  p_i(T) \;=\;
  \frac{\exp(z_i / T)}{\sum_j \exp(z_j / T)},
  \label{eq:temperatura}
\end{equation}
em que $z_i$ são os \emph{logits}; com $T > 1$, as probabilidades pequenas
deixam de ser esmagadas. O aluno é treinado numa combinação de duas perdas:
a entropia cruzada usual contra os rótulos verdadeiros e uma perda de
imitação contra os \emph{soft targets} do professor na mesma temperatura.
Em regressão, a receita análoga substitui o rótulo pelo valor predito pelo
professor --- e, na extensão distribucional do Capítulo~7, pela \emph{curva
de quantis} inteira do professor.

\paragraph{Destilar para árvores.} A linhagem \emph{born-again} mostrou
cedo que o aluno não precisa ser uma rede: \citet{breiman1996} comprimem um
\emph{ensemble} de árvores numa única árvore treinada nas predições do
comitê, com perda de fidelidade surpreendentemente pequena;
\citet{vidal2020} dão à compressão de \emph{ensembles} tratamento exato
moderno. Na direção inversa, \citet{frosst2017} destilam uma rede profunda
numa árvore de decisão suave em busca de interpretabilidade, e a destilação
de redes em redes de mesma capacidade (\emph{born-again networks}) revelou
que até a auto-imitação pode melhorar o aluno \textbf{[REF-NOVA:
Furlanello et al. 2018, \emph{Born-Again Neural Networks}]}. A lição
transversal: a destilação é agnóstica ao par professor--aluno; o que se
transfere é a \emph{função} aprendida, exposta pelas saídas.

\paragraph{Destilar \emph{foundation models} tabulares.} Com professores
ICL, a destilação ganha um atrativo adicional: ela remove exatamente o
gargalo constitutivo da família (o contexto que viaja a cada predição),
convertendo um professor de inferência cara num aluno de inferência
barata. O \emph{Pocket FM} \citep{pocketfm2026} estabelece essa agenda para
\emph{classificação}, em escala, com uma salvaguarda metodológica que o
Capítulo~7 herda: a rotulagem \emph{out-of-fold} --- o professor rotula
cada exemplo condicionado apenas nas demais partições, nunca nele próprio
--- para que o aluno não aprenda dos vazamentos do professor. O que
permanece em aberto nessa linha, e constitui a contribuição desta
dissertação, é a extensão à \emph{regressão distribucional}: destilar não a
média, mas a distribuição preditiva --- o que exige o ferramental da
próxima seção. O Capítulo~3 situa esse recorte no estado da arte.

\section{Regressão distribucional: quantis, CRPS e calibração}
\label{sec:fund-distribucional}

Prever um número é útil; prever um número \emph{com a incerteza em volta} é
o que decisões de risco exigem. Considere um modelo que avalia imóveis:
dizer ``este imóvel vale 500 mil'' e dizer ``vale 500 mil, e com 80\% de
probabilidade entre 430 e 580 mil'' são produtos diferentes --- o segundo
permite precificar risco, definir margens e recusar decisões quando a
incerteza é grande demais. Dois modelos com o mesmo erro pontual médio
podem diferir radicalmente na qualidade dessa segunda resposta, e as
métricas pontuais (RMSE, MAE) são cegas à diferença. Esta seção apresenta
o instrumental que a enxerga --- quantis e perda \emph{pinball}, CRPS,
calibração e afiação --- na profundidade que o Capítulo~7 pressupõe: aqui
a matemática é detalhada de propósito, pois estas são as grandezas da
contribuição.

\subsection{Quantis e a perda \emph{pinball}}

O quantil condicional $q_\tau(\mathbf{x})$ é o valor abaixo do qual o alvo
cai com probabilidade $\tau$, dado $\mathbf{x}$: o quantil $\tau = 0{,}5$ é
a mediana condicional; os quantis $0{,}1$ e $0{,}9$ delimitam um intervalo
central de 80\%. A regressão de quantis \citep{koenker1978} estima
$q_\tau$ diretamente, minimizando a \emph{perda pinball} (ou perda de
quantil): com erro $u = y - \hat{q}_\tau$,
\begin{equation}
  \rho_\tau(u) \;=\;
  \begin{cases}
    \tau\, u, & u \geq 0 \quad \text{(subestimou o alvo)},\\[2pt]
    (\tau - 1)\, u, & u < 0 \quad \text{(superestimou o alvo)}.
  \end{cases}
  \label{eq:pinball}
\end{equation}

\paragraph{A intuição assimétrica.} A perda é uma rampa com duas
inclinações (Figura~\ref{fig:pinball}): para $\tau = 0{,}9$, subestimar
custa $0{,}9$ por unidade de erro e superestimar custa apenas $0{,}1$. O
minimizador ótimo, portanto, prefere errar para cima --- e erra para cima
na medida exata: a predição que minimiza a perda esperada é precisamente o
quantil $0{,}9$ da distribuição condicional. É essa propriedade --- o
minimizador da perda \emph{pinball} de nível $\tau$ é o quantil $\tau$ ---
que faz dela o instrumento canônico da regressão de quantis; com
$\tau = 0{,}5$, ela se reduz (a um fator) ao erro absoluto, cujo minimizador
é a mediana. Treinar um modelo numa \emph{grade} de níveis
$\tau \in \{0{,}05, \ldots, 0{,}95\}$ produz uma curva de quantis ---
uma discretização da distribuição preditiva inteira.

\begin{figure}[htb]
\centering
\begin{tikzpicture}[font=\small, line cap=round]
  % eixos
  \draw[-{Stealth[length=2mm]}] (-2.6,0) -- (2.9,0)
    node[below right=0pt] {$u = y - \hat{q}_\tau$};
  \draw[-{Stealth[length=2mm]}] (0,-0.25) -- (0,2.1)
    node[above] {$\rho_\tau(u)$};
  \node[below, font=\footnotesize] at (-1.3,0) {superestimação};
  \node[below, font=\footnotesize] at (1.45,0) {subestimação};
  % tau = 0.5 (simétrica, tracejada)
  \draw[dashed, gray, thick] (-2.3,1.15) -- (0,0) -- (2.3,1.15);
  \node[gray, font=\footnotesize, anchor=south east] at (-1.6,0.9)
    {$\tau = 0{,}5$};
  % tau = 0.9 (assimétrica, sólida)
  \draw[very thick] (-2.3,0.23) -- (0,0) -- (2.3,2.07);
  \node[font=\footnotesize, anchor=south east] at (2.35,1.5)
    {$\tau = 0{,}9$};
\end{tikzpicture}
\caption{\textbf{A perda \emph{pinball} penaliza os dois lados do erro com
inclinações diferentes, e é essa assimetria que aponta o minimizador para o
quantil~$\tau$.} Para $\tau = 0{,}9$ (linha sólida), subestimar o alvo
custa nove vezes mais que superestimar, de modo que a predição ótima se
posiciona no ponto em que o alvo só a excede 10\% das vezes --- o quantil
$0{,}9$. Em $\tau = 0{,}5$ (tracejada), a perda é simétrica e recupera o
erro absoluto, cujo minimizador é a mediana.}
\label{fig:pinball}
\end{figure}

\paragraph{Um efeito colateral da grade: o cruzamento de quantis.} Quando
os níveis da grade são estimados de forma independente, nada garante a
monotonicidade que a definição exige --- o modelo pode produzir
$\hat{q}_{0{,}8} < \hat{q}_{0{,}7}$ para alguma linha, o chamado
\emph{cruzamento de quantis}, que torna a ``distribuição'' internamente
inconsistente. O reparo pragmático padrão é reordenar os quantis preditos
de cada linha em ordem crescente, operação que nunca piora a perda
\emph{pinball} total e restaura a validade da curva; é o reparo adotado
nos alunos multi-quantil do Capítulo~7.

\subsection{CRPS: uma régua para distribuições}

Como comparar duas distribuições preditivas? A régua desta dissertação é o
\emph{continuous ranked probability score} (CRPS): para uma distribuição
preditiva com função de distribuição acumulada $F$ e um alvo observado $y$,
\begin{equation}
  \operatorname{CRPS}(F, y) \;=\;
  \int_{-\infty}^{\infty}
  \left( F(z) - \mathbb{1}\{z \geq y\} \right)^2 \mathrm{d}z,
  \label{eq:crps}
\end{equation}
a distância quadrática integrada entre a CDF predita e a CDF degrau do
valor observado. Duas propriedades o qualificam. Primeiro, o CRPS é uma
\emph{regra de pontuação própria} (\emph{proper scoring rule})
\citep{gneiting2007}: seu valor esperado é minimizado quando o modelo
reporta a verdadeira distribuição condicional --- não há como ganhar
pontuação ``blefando'' uma distribuição diferente da que o modelo de fato
acredita, o que o torna o critério honesto para avaliar incerteza.
Segundo, ele \emph{generaliza o erro absoluto}: se $F$ colapsa num ponto
(predição pontual), a Eq.~\ref{eq:crps} reduz-se a $|{\hat{y} - y}|$ --- de
modo que o CRPS permite comparar, na mesma régua e na unidade do alvo,
modelos pontuais e distribucionais. Menor é melhor.

\paragraph{O estimador de grade.} Na prática, quando o modelo entrega uma
grade de quantis em vez de uma CDF contínua, o CRPS é estimado como duas
vezes a média da perda \emph{pinball} sobre a grade de níveis
\citep{gneiting2011},
\begin{equation}
  \widehat{\operatorname{CRPS}}(F, y) \;=\;
  \frac{2}{|\mathcal{T}|} \sum_{\tau \in \mathcal{T}}
  \rho_\tau\!\left(y - \hat{q}_\tau\right),
  \label{eq:crps-grade}
\end{equation}
o estimador padrão da prática de previsão probabilística --- e o elo formal
entre as duas grandezas desta seção: minimizar a perda \emph{pinball} em
toda a grade é, a uma constante, otimizar o CRPS discretizado. Por ser uma
discretização, o estimador carrega um viés dependente da grade; o
Capítulo~7 reporta a validação desse viés no protocolo desta dissertação e
a regra de leitura que dele decorre (comparações internas, nunca valores
absolutos entre trabalhos).

\subsection{Calibração e afiação}

O CRPS resume a qualidade distribucional num número; os diagnósticos de
\emph{calibração} a decompõem. Um intervalo preditivo nominal de 80\% está
\emph{calibrado} se cobre o alvo em cerca de 80\% dos casos de teste. A
métrica operacional é o PICP (\emph{prediction interval coverage
probability}): a fração dos alvos de teste que caem dentro do intervalo
nominal --- PICP80 muito abaixo de 80\% denuncia um modelo
\emph{superconfiante} (intervalos estreitos demais), muito acima, um modelo
\emph{subconfiante} (intervalos largos demais). Cobertura, porém, é
condição necessária e não suficiente: um intervalo absurdamente largo cobre
sempre e não informa nada. Por isso a cobertura é lida junto com a
\emph{afiação} (\emph{sharpness}) --- a largura média dos intervalos: entre
modelos igualmente calibrados, o de intervalos mais estreitos é o mais
informativo \textbf{[REF-NOVA: Gneiting, Balabdaoui \& Raftery 2007,
paradigma calibração--afiação]}.

\subsection{Dois \emph{baselines} distribucionais nativos}

Fecham o instrumental dois modelos clássicos que produzem distribuições
preditivas sem professor algum --- os controles naturais de qualquer
proposta de destilação distribucional. As \emph{quantile regression
forests} (QRF) estendem as florestas aleatórias guardando, em cada folha,
não a média mas \emph{todos} os valores de treino que ali caem; a
distribuição preditiva de uma linha nova é a distribuição empírica
ponderada desses valores, da qual se lê qualquer quantil
\textbf{[REF-NOVA: Meinshausen 2006, \emph{Quantile Regression Forests}]}.
O \emph{NGBoost} segue a rota paramétrica: um \emph{gradient boosting} cujas
saídas são os \emph{parâmetros} de uma distribuição (por exemplo, média e
variância de uma gaussiana), treinado por gradientes naturais sobre uma
regra de pontuação própria \textbf{[REF-NOVA: Duan et al. 2020,
\emph{NGBoost}]}. QRF representa a via não-paramétrica; NGBoost, a
paramétrica --- e ambos aparecem como competidores no Capítulo~7.

\section{Comparação estatística de algoritmos}
\label{sec:fund-estatistica}

Resta o último instrumento: como afirmar, com honestidade estatística, que
um modelo é melhor que outro --- ou que dois modelos empatam --- a partir
de resultados em uma coleção de conjuntos de dados? Esta seção apresenta a
teoria dos testes que o Capítulo~4 adota como protocolo e o Capítulo~5
aplica; a exposição segue o arcabouço canônico da área
\citep{demsar2006,benavoli2017}.

\paragraph{Por que não uma ANOVA sobre as métricas?} Porque as métricas não
são \emph{comensuráveis} entre conjuntos de dados: uma AUC de $0{,}95$ num
conjunto fácil e uma de $0{,}75$ num difícil não estão na mesma escala, e
médias ou variâncias entre conjuntos misturam alhos com bugalhos --- além
de a normalidade, pressuposto da ANOVA, não ter por que valer. A solução de
\citet{demsar2006} é trabalhar com \emph{ranks}: em cada conjunto de dados,
ordenam-se os $k$ modelos do melhor ($1$) ao pior ($k$); os ranks, ao
contrário das métricas, são comparáveis entre conjuntos por construção.

\paragraph{Qual é a unidade de análise?} Uma segunda decisão precede
qualquer teste: o que conta como uma observação independente? Quando cada
modelo é avaliado por validação cruzada, é tentador tratar cada
\emph{fold} como uma réplica --- mas os folds de um mesmo conjunto de
dados compartilham a maior parte dos exemplos e não são independentes;
usá-los como réplicas inflaria artificialmente o $N$ efetivo e, com ele, o
poder aparente dos testes (\emph{pseudo-replicação}). A unidade correta é
o \emph{conjunto de dados}: cada um contribui com um único número por
modelo (a média entre folds), e os folds alimentam apenas a descrição da
variabilidade intra-conjunto. Todos os testes a seguir pressupõem essa
escolha.

\paragraph{Friedman: o teste \emph{omnibus}.} O teste de Friedman pergunta
se as diferenças entre os ranks médios $R_j$ dos $k$ modelos sobre os $N$
conjuntos são explicáveis pelo acaso \textbf{[REF-NOVA: Friedman 1937,
teste de ranks de medidas repetidas]}. Sob a hipótese nula de equivalência
entre todos os modelos, a estatística
\begin{equation}
  \chi^2_F \;=\; \frac{12N}{k(k+1)}
  \left[ \sum_{j=1}^{k} R_j^2 \;-\; \frac{k(k+1)^2}{4} \right]
  \label{eq:friedman}
\end{equation}
segue aproximadamente uma distribuição qui-quadrado com $k-1$ graus de
liberdade. Essa aproximação é conservadora quando $N$ é pequeno --- o
regime desta dissertação ---, e por isso reporta-se também a correção de
Iman--Davenport,
\begin{equation}
  F_F \;=\; \frac{(N-1)\,\chi^2_F}{N(k-1) - \chi^2_F},
  \label{eq:iman-davenport}
\end{equation}
que segue uma distribuição $F$ e rejeita com mais poder \textbf{[REF-NOVA:
Iman \& Davenport 1980, correção do teste de Friedman]}. Um Friedman
significativo diz apenas que \emph{alguma} diferença existe; localizá-la é
tarefa dos testes \emph{post-hoc}.

\paragraph{Nemenyi e o diagrama de diferença crítica.} O \emph{post-hoc} de
Nemenyi compara todos os pares de modelos pelos ranks médios, com controle
embutido do erro de família \textbf{[REF-NOVA: Nemenyi 1963, comparações
múltiplas por ranks]}: dois modelos diferem significativamente se seus
ranks médios distam mais que a \emph{diferença crítica}
\begin{equation}
  \mathrm{CD} \;=\; q_\alpha \sqrt{\frac{k(k+1)}{6N}},
  \label{eq:cd}
\end{equation}
em que $q_\alpha$ vem da distribuição da amplitude studentizada
\citep{demsar2006}. O resultado é usualmente apresentado num \emph{diagrama
de diferença crítica} (diagrama CD), e vale a pena aprender a lê-lo, pois o
Capítulo~5 o utiliza repetidamente: os modelos são dispostos sobre um eixo
horizontal na posição de seu rank médio (melhores de um lado, piores do
outro); uma régua indica o comprimento do CD; e barras horizontais espessas
conectam grupos de modelos cujos ranks médios distam menos que o CD ---
\emph{modelos ligados pela mesma barra não são estatisticamente
distinguíveis} àquele nível. Duas leituras erradas a evitar: a barra não
diz que os modelos conectados são iguais (apenas que o teste não os
distingue --- ausência de evidência não é evidência de ausência), e a
Eq.~\ref{eq:cd} mostra que o CD cresce com $k$ e encolhe lentamente com
$N$: com muitos modelos e poucos conjuntos, barras largas são a
consequência aritmética esperada, não uma anomalia.

\paragraph{Wilcoxon pareado com correção de Holm.} Para pares específicos
de modelos, o teste de postos sinalizados de Wilcoxon é mais poderoso que o
Nemenyi, porque usa as \emph{magnitudes} das diferenças por conjunto de
dados, não apenas os ranks globais \textbf{[REF-NOVA: Wilcoxon 1945, teste
de postos sinalizados]}. Ao testar muitos pares, porém, multiplica-se a
chance de falsos positivos; o procedimento de Holm corrige isso ordenando
os p-valores e exigindo limiares progressivamente mais frouxos ---
controle do erro de família menos conservador que a correção de Bonferroni
simples \textbf{[REF-NOVA: Holm 1979, procedimento sequencial de rejeição]}.
Uma limitação aritmética importa no regime desta dissertação: com $N$
conjuntos de dados, o menor p-valor bilateral que o Wilcoxon pode produzir
é $2 \cdot (1/2^N)$ --- com $N = 5$, o piso é $0{,}0625$, e rejeitar a
$0{,}05$ é \emph{matematicamente impossível}, antes de qualquer correção. O
Capítulo~4 declara as consequências protocolares desse piso.

\paragraph{O que significa ``praticamente equivalente''? O teste bayesiano
com ROPE.} O arcabouço frequentista tem uma assimetria incômoda: ele pode
rejeitar a igualdade, mas não pode \emph{afirmá-la} --- não-rejeição não é
evidência de equivalência, especialmente com $N$ pequeno e pouco poder.
\citet{benavoli2017} oferecem o complemento bayesiano: o teste de postos
sinalizados bayesiano estima a distribuição \emph{posterior} da diferença
média de desempenho entre dois modelos e introduz uma \emph{região de
equivalência prática} (ROPE, \emph{region of practical equivalence}) --- um
intervalo em torno de zero, fixado \emph{a priori} por juízo de domínio,
dentro do qual diferenças são irrelevantes na prática. A posterior é então
resumida em três probabilidades: a de o primeiro modelo ser praticamente
melhor, a de o segundo ser praticamente melhor e a de os dois serem
\emph{praticamente equivalentes} (a massa da posterior dentro da ROPE).
Dizer que dois modelos são praticamente equivalentes é, portanto, uma
afirmação positiva e quantificada --- ``com probabilidade $p$, a diferença
entre eles é menor do que o que declaramos irrelevante'' --- e não um
não-resultado. É esse instrumento que autoriza as conclusões de empate do
Capítulo~5; a escolha do raio da ROPE e sua análise de sensibilidade são
decisões de protocolo, especificadas no Capítulo~4.

\paragraph{Tamanho de efeito e decomposição de variância.} Significância
não mede relevância, e dois complementos fecham o arcabouço. O tamanho de
efeito de Cohen na variante pareada,
$d_z = \bar{d} / s_d$ (média das diferenças por conjunto dividida pelo seu
desvio-padrão), quantifica \emph{quanto} dois modelos diferem na escala da
própria variabilidade das diferenças. E a razão de correlação
$\eta^2 = \mathit{SQ}_{\text{entre}} / \mathit{SQ}_{\text{total}}$ ---
aqui aplicada aos ranks, com os modelos agrupados por família ---
decompõe a variância total dos ranks na parcela explicada pelo
agrupamento: $\eta^2$ alto diria que a família (GBDT, aprendizado profundo,
\emph{foundation model}) determina o desempenho; $\eta^2$ baixo, que a
variação dentro das famílias domina e a escolha do modelo individual
importa mais que a do paradigma. O Capítulo~5 usa exatamente essa leitura.

\begin{transicao}
Este capítulo montou o vocabulário: o objeto (dados tabulares), os
competidores (GBDTs, aprendizado profundo, \emph{foundation models} com
ICL), as técnicas transversais (destilação, regressão distribucional) e a
régua (testes estatísticos sobre ranks, com a camada bayesiana de
equivalência). O próximo capítulo usa esse vocabulário para mapear o estado
da arte --- os benchmarks tabulares existentes e seus protocolos, a
literatura de destilação de \emph{foundation models} tabulares e as
avaliações distribucionais independentes --- e para demarcar, com datas e
verificações explícitas, as lacunas que as três contribuições desta
dissertação ocupam.
\end{transicao}
